\section{Finite quotients in odd division degree}\label{sec:quotients}

Assume throughout this section that $D$ is a division algebra of odd
reduced degree $n\geq3$ over the quadratic extension $L/k$.  Put
$S=\SU(D,*)$ and $U=\U(D,*)$, as in Section~\ref{sec:arithmetic}.
The one-dimensional Hermitian space over $D$ is anisotropic, so $S$
is $k$-anisotropic.  The reduced norm gives a surjective homomorphism
\[
 \det:U(k)\longrightarrow L^1,
 \qquad L^1=\{a\in L^\times:a\bar a=1\},
\]
with kernel $S(k)$.  Recall that, for a fixed integer $e\geq1$,
\[
 R=S(k)^e,\qquad
 P_0=\overline{\delta_A(R)},\qquad
 V_0=\delta_A^{-1}(P_0),\qquad V=V_0/R.
\]
Here $S(k)^e$ denotes the subgroup generated by the $e$th powers.
The group $V$ is finite by Proposition~\ref{arith:power-defect}.
Our purpose is to turn a nontrivial quotient of $V$ into a finite
quotient defined on the full unitary group.

\begin{proposition}[Finite quotient dichotomy]\label{quo:dichotomy}
If $V\ne1$, at least one of the following homomorphisms exists:
\begin{enumerate}
 \item a homomorphism $\chi:U(k)\to B$ to a finite abelian group such
 that $\chi(S(k))\ne1$;
 \item a homomorphism $\phi:U(k)\to\mathcal F$ to a finite nonabelian
 simple group such that
 \[
   \phi(V_0)=\mathcal F,\qquad R\subseteq\ker\phi.
 \]
\end{enumerate}
\end{proposition}

The argument follows the quotient construction of
\cite[Section~3, especially Propositions~3.1 and~3.2]{MPNF}.
We give the details because two parts require care in positive
characteristic: commuting elements must be constructed in separable
field tori, and the treatment of local principal units must not assume
that the subgroup of $n$th powers is open.  We first prove the commuting
approximation statement that controls the conjugation action on $V$.

\subsection{Commuting elements with prescribed determinants}

At a place $v\in A$, the extension $L/k$ splits and
\[
 U(k_v)=\mathcal D_v^\times,
 \qquad S(k_v)=\SL_1(\mathcal D_v),
\]
where $\mathcal D_v$ is a central division algebra of degree $n$ over
$k_v$.  A \emph{separable field torus} in $U(k_v)$ means the
multiplicative group of a maximal separable subfield of $\mathcal D_v$.
This qualification matters when the characteristic divides $n$.

\begin{proposition}[Commuting approximation]\label{quo:commuting}
The following assertions hold.
\begin{enumerate}
 \item Let $a,b\in L^1$.  For each $v\in A$, let $x_v,t_v\in U(k_v)$
 belong to one common separable field torus and satisfy
 $\det x_v=a_v$ and $\det t_v=b_v$.
 There exist commuting elements $x,t\in U(k)$ with
 \[
  \det x=a,\qquad \det t=b,
 \]
 whose components at $A$ approximate the prescribed pairs arbitrarily
 closely.
 \item Given $\gamma\in V_0$ and $h\in U(k)$, there exist commuting
 elements $x,t\in U(k)$ such that
 \[
  x\in\gamma R,\qquad \det t=\det h,
 \]
 and $t_A$ is arbitrarily close to $h_A$.
\end{enumerate}
\end{proposition}

\begin{proof}
In the second assertion set $a=1$ and $b=\det h$.  In both assertions
we shall construct a regular semisimple first element $x$ whose
centralizer has local elements of determinant $b$ at every place.
The norm-torus theorem of Section~\ref{sec:norm-tori} will then give
a global second element, together with its approximations at $A$.

We begin with the local conditions at $A$.  In the first assertion,
perturb $x_v$ within the determinant-$a_v$ fiber of its given field
torus until it is regular semisimple.  The field is separable, so its
norm is a smooth map.  After extension to an algebraic closure, a norm
fiber is a translate of
\[
 \{(z_1,\ldots,z_n)\in\bbG_m^n:z_1\cdots z_n=1\}.
\]
No equality $z_i=z_j$, with $i\ne j$, holds identically on this fiber.
Its regular elements therefore form a dense open subset, also dense
in the local topology.  Near a regular semisimple element the
conjugation map from the group times its torus is a submersion.
Consequently the centralizers of nearby elements are obtained by small
conjugations.  The same conjugations transport $t_v$, preserving its
determinant and the required approximation.

For the second assertion, first perturb $h_v$ to a regular semisimple
element in its determinant fiber.  Such a perturbation is possible
because the fiber is a smooth translate of $S$ and its regular
semisimple open is nonempty.  In the resulting separable field torus
choose a regular norm-one element close to $1$ as the first element.
The preceding conjugation argument again gives a local opening for
$x$.  These openings near $1$ are compatible with the prescribed coset
$\gamma R$: its closure at $A$ is $P_0$, and $P_0$ is open.

Choose a finite set $B$ containing $A$, an isotropic place to be omitted
in strong approximation, all exceptions for integral models, all places
ramified in $L/k$, and the places where $a$ or $b$ is not a unit.
At the additional places we need local tori whose determinant images
contain both $a_v$ and $b_v$.  At a nonsplit place of $L/k$, the algebra
splits and a diagonal Hermitian torus has determinant image equal to
the local norm-one group.  At a split place outside $A$, write a
component of the algebra as $M_l(\Delta)$.  If $\Delta$ is nontrivial,
then $l\geq2$.  Use an unramified maximal subfield of $\Delta$ in one
diagonal block and a totally ramified separable maximal subfield in a
second block.  The first norm supplies all units, and the second norm
supplies an element of valuation one; their combined image is all of
$k_v^\times$.  Such separable totally ramified fields exist in every
characteristic, for instance by separable Eisenstein polynomials, and
embed in $\Delta$ by the local invariant criterion.  Complete the
remaining blocks to a maximal \'etale subalgebra.  If $\Delta=k_v$, use
the ordinary diagonal torus.  In every case choose the first element
regular with determinant $a_v$ and take a sufficiently small opening
around it.  These constructions also give nonempty local openings at
any further exceptional places introduced by the approximation theorem.

We add finitely many auxiliary places, split in $L$ and in $D$, at
which $a$ and $b$ are units.  Prescribe unramified regular centralizers
having every possible permutation cycle type in $S_n$, one type at
each place.  Unramified norms supply both unit determinants.  We may
use places with sufficiently large residue fields that regular elements
with those determinants exist.  These conditions will force the Galois
group required by Proposition~\ref{tor:norm-approximation}.

We now verify the hypotheses of Proposition~\ref{approx:power-factor}
on the determinant-$a$ slice $X$.  This slice is a rational translate
of $S$, since the determinant on $U(k)$ is onto.  For the first
assertion, weak approximation supplies a rational point meeting the
finitely many openings; we then apply the proposition to its coset
modulo $R$.  For the second assertion, we apply it to $\gamma R$.
The closure statement of Proposition~\ref{arith:closure} permits the
openings near $1$ at $A$ and places no restriction on the other local
openings.

Define $Y\subset X$ over $k$ to be the closure of the geometric locus
whose reduced characteristic polynomial has no root of multiplicity one.
The root-multiplicity condition is preserved by the descent defining
$X$.  We shall require integral reductions outside $B$ to avoid the
spread-out model of $Y$.  It has geometric codimension at least two
in $X$.  Indeed, on a split matrix component there are at most $\lfloor n/2\rfloor$ distinct eigenvalues on this locus.  The
fixed determinant reduces the eigenvalue parameter dimension by one,
even if the resulting torus equation is inseparable.  A conjugacy
class with fixed eigenvalues has dimension at most $n^2-n$.
There are finitely many Jordan types, so
\[
 \dim Y\leq n^2-n+\lfloor n/2\rfloor-1
             \leq (n^2-1)-2.
\]
This codimension bound and the description by root multiplicities
persist outside finitely many model exceptions.

Suppose that an integral regular semisimple element reduces outside
$Y$, and write $K=k_v$ and $L_K=L\otimes_k K$.  A simple absolute
residue root belongs to an irreducible residue factor of multiplicity
one.  Hensel lifting supplies an unramified field factor of the \'etale
centralizer over $L_K$.  Take its full orbit under the involution and
denote the resulting algebra by $E_j$.  All its factors are unramified
over $K$, since $L_K/K$ is unramified or split.  If the field factor is
stable under the involution, its fixed field $F_j$ is unramified over
$K$.  If two factors are exchanged, their fixed algebra is the diagonal
field under that exchange, and is again unramified.  In both cases
quadratic descent identifies the algebra with involution as
\[
 E_j=F_j\otimes_K L_K,
\]
with the identity on $F_j$ and quadratic conjugation on $L_K$.

Its unitary torus contributes local determinant norms as follows.
After an unramified extension splitting these algebras, it has one
free multiplicative coordinate $z_\tau$ for each embedding
$\tau:F_j\hookrightarrow\overline K$.  The determinant is
$\prod_\tau z_\tau^{\epsilon_\tau}$ with
$\epsilon_\tau\in\{1,-1\}$: each field embedding occurs once in the
reduced characteristic representation, and the sign records which
member of its involution pair lies in the chosen $L_K$-component.
Passing to the other member inverts the coordinate.  The dual map from the rank-one
character lattice of $L_K^1$ therefore has primitive image.  Thus the
determinant is a surjection of unramified tori with a torus as kernel,
in every characteristic.  Lang's theorem makes the map onto on
residue-field points, and smooth lifting makes it onto integral
points.  It supplies the unit $b_v$.

For a single pole, the approximation theorem places us at a place split
in the fixed Galois field and gives
\[
 x=g_L\diag(t^{e d_1},\ldots,t^{e d_n})g_R,
 \qquad d_1<\cdots<d_n,
\]
with $\ord_v(t)>0$ and integral invertible side factors.
Require the principal minor on the first $j$ indices of $g_Rg_L$ to
be a unit for each $1\leq j\leq n$.  These are nonempty algebraic open
conditions on each pole hyperplane.  Indeed, the cyclic product
$g_Rg_L$ contains one entire undamaged basic parameter list.  Keeping
all other parameters fixed, that list varies $g_Rg_L$ densely in the
fixed determinant slice.  The principal-minor conditions are nonempty
on this slice, as witnessed by a diagonal matrix of the given
determinant.  Consequently, once the two basic lists have been chosen,
every initial $x$ and every appended conjugator tuple admit these
pole tests; no further generic restriction on those choices is needed.
Identity values of appended parameters preserve the original tests.
If $c_j$ is the $j$th elementary characteristic coefficient, its principal-minor expansion gives
\begin{equation}\label{quo:pole-slopes}
 \ord_v(c_j)=e(d_1+\cdots+d_j)\ord_v(t).
\end{equation}
The selected minor is the unique term of least valuation.  Thus all
Newton segments have length one and distinct slopes.  The
characteristic polynomial splits into distinct linear factors over
$k_v$, so its centralizer supplies every determinant.

For the allowed sets in Proposition~\ref{approx:power-factor}, include
all elements outside the regular semisimple open.  On that open use
the determinant-norm condition just proved.  The integral and pole
conditions hold, and the local tori constructed above provide the
required nonempty openings.  Impose regularity at one place of $B$.
The proposition now gives a global regular semisimple $x$, with the
prescribed determinant, local conditions, and, in the second assertion,
its exact coset modulo $R$.

Since $D$ is division, $E=L(x)$ is a maximal separable field in $D$.
It is stable under $*$ because $x^*=x^{-1}$.  If $F=E^{*=1}$ is its
fixed field, then $E=LF$ and $[F:k]=n$.  Let $\Gamma$ be the Galois
group of the joint normal closure of $F$ and $L$; its natural action
embeds it in $S_n\times C_2$.  The prescribed Frobenius classes at
places split in $L$ imply that the kernel of $\Gamma\to C_2$ meets
every conjugacy class of $S_n$.  Such a subgroup is all of $S_n$.
For, if it were proper, the transitive action on its cosets would
have a derangement: the average number of fixed cosets is one, while
the identity fixes more than one.  A derangement cannot be conjugate
to an element of the subgroup.  Since $\Gamma$ also surjects onto
$C_2$, we obtain $\Gamma=S_n\times C_2$.

The unitary centralizer of $x$ is
\[
 T_E=\Res_{F/k}\bigl(\Res^1_{E/F}\bbG_m\bigr),
 \qquad \det|_{T_E}=N_{E/L}.
\]
The fiber $N_{E/L}^{-1}(b)$ has points at every completion, by the
local openings, integral conditions, and pole conditions.  It is a
torsor under the norm kernel of Proposition~\ref{tor:norm-approximation},
whose Galois hypothesis we have just checked.  That theorem gives a
rational point $t$ with the desired approximations at $A$.
It lies in the centralizer of $x$ and has determinant $b$.
\end{proof}

\subsection{The action of the full unitary group}

The group-theoretic passage from a class-preserving action to a quotient
of the ambient group is adapted from
\cite[Corollary~2.3, Theorem~2.4, and the proof of Theorem~2.1]{RapinchukPotapchik1996}.
Here Proposition~\ref{quo:commuting} supplies the unitary input.
Conjugation by $U(k)$ preserves $R$, which is characteristic in
$S(k)$, and consequently preserves $P_0$ and $V_0$.  We show that its
action on $V$ preserves every conjugacy class.  Take $\gamma\in V_0$
and $h\in U(k)$.  Proposition~\ref{quo:commuting}(ii) gives
$x\in\gamma R$ commuting with $t$, where $\det t=\det h$.
Choose the approximation at $A$ sufficiently close that
$th^{-1}\in V_0$.  Conjugation by $t$ fixes $\gamma R$, while
conjugation by $th^{-1}$ is inner on $V$.  Hence conjugation by $h$
sends $\gamma R$ to a conjugate of itself.

Suppose now that $V\ne1$, and choose a simple quotient $C$ of $V$.
Every normal subgroup of $V$ is a union of conjugacy classes, so the
kernel of $V\to C$ is preserved by $U(k)$.  The induced automorphisms
of $C$ are class-preserving.  If $C$ is nonabelian, the theorem of
Feit--Seitz \cite[Theorem~C]{FeitSeitz} makes each of these automorphisms
inner.  As $C$ has trivial center, the action gives a homomorphism
\[
 U(k)\longrightarrow\Inn(C)=C.
\]
On $V_0$ it is the original quotient map and it kills $R$.
This proves the second alternative of Proposition~\ref{quo:dichotomy}.

It remains to handle the case in which $C$ is cyclic of prime order.
Let $R'$ be its kernel upstairs.  Then
\begin{equation}\label{quo:central-quotient}
 R\subseteq R'\subseteq V_0,
 \qquad C=V_0/R'\subseteq Z\bigl(U(k)/R'\bigr).
\end{equation}
The rest of the section uses this nonzero central quotient to construct
a finite character of $U(k)$ that is nonzero on $S(k)$.  We shall use the compact
local groups at $A$ to construct and split a circle extension.

\subsection{Characters of a local norm-one group}

The required local calculation is the division-algebra commutator
calculation underlying \cite{Riehm}; we include its filtration proof,
using the standard cyclic description of local division algebras
\cite{Reiner}.  It explains why divisibility of $n$ by the
characteristic causes no additional abelian quotient.

\begin{lemma}\label{quo:local-abelianization}
Let $K$ be a nonarchimedean local field of positive characteristic,
with residue field $\bbF_q$, and let $\mathcal D$ be a central division
algebra of odd degree $n\geq3$ over $K$.  Set $S_K=\SL_1(\mathcal D)$
and
\[
 T_K=\ker\bigl(N_{\bbF_{q^n}/\bbF_q}:
                    \bbF_{q^n}^{\times}\to\bbF_q^{\times}\bigr).
\]
Then reduction induces
\[
 S_K/\overline{[S_K,S_K]}\simeq T_K,
 \qquad
 \overline{[\mathcal D^\times,\mathcal D^\times]}=S_K.
\]
There is a generator $\sigma$ of
$\Gal(\bbF_{q^n}/\bbF_q)$ such that conjugation by an element of
reduced-norm valuation $r$ acts on $T_K$ as $\sigma^r$.
\end{lemma}

\begin{proof}
Choose an unramified maximal subfield $E_K/K$ and a prime element
$\Pi$ of $\mathcal D$ inducing a generator $\sigma$ on $E_K$.
We normalize so that $\ord_K(\Nrd\Pi)=1$; $\Pi^n$ is a central
uniformizer.  Let $\mathfrak P=(\Pi)$ be the maximal ideal of the
maximal order of $\mathcal D$, and let $\mathfrak p$ be the maximal
ideal of $K$.  Every element of $S_K$ is integral.  Its reduction has
norm one, and the norm-one Teichm\"uller representatives in $E_K$
give a multiplicative section of the reduction map onto $T_K$.

Put $S_r=S_K\cap(1+\mathfrak P^r)$ for $r\geq1$.  The norm filtration
is
\begin{equation}\label{quo:norm-filtration}
 \Nrd(1+\mathfrak P^r)=1+\mathfrak p^{\lceil r/n\rceil}.
\end{equation}
For the inclusion, the roots of the reduced characteristic polynomial
of an element of $\mathfrak P^r$ have valuation at least $r/n$;
their elementary symmetric coefficients have the corresponding
valuation bounds.  For surjectivity, use norms of principal units in
the unramified field $E_K$: its norm maps
$1+\mathfrak p_{E_K}^{\lceil r/n\rceil}$ onto the right-hand side.
The latter assertion follows successively from surjectivity of the
finite-field trace and completeness.

Leading-coefficient coordinates now give
\begin{equation}\label{quo:norm-one-grades}
 S_r/S_{r+1}\simeq
 \begin{cases}
  (\bbF_{q^n},+),&n\nmid r,\\
  \ker\Tr_{\bbF_{q^n}/\bbF_q},&n\mid r.
 \end{cases}
\end{equation}
If $n\nmid r$, the norm images at levels $r$ and $r+1$ agree in
\eqref{quo:norm-filtration}, so any leading coefficient can be
corrected at the next level to give norm one.  If $r=nm$, the first
norm coefficient of $1+a\Pi^r$ is the residue trace of $a$ at
level $m$; correction at the next level is possible exactly when
that trace is zero.  The trace map of a finite-field extension is
surjective even when the characteristic divides $n$.

For $n\nmid r$, choose a Teichm\"uller element $t\in T_K$ not fixed
by $\sigma^r$.  Such a $t$ exists because
$|T_K|=(q^n-1)/(q-1)$ is larger than the multiplicative group of any
proper subfield of $\bbF_{q^n}$.  The commutator with $t$ multiplies
the leading coefficient at level $r$ by the nonzero scalar
$t/\sigma^r(t)-1$.  It therefore fills the grade in
\eqref{quo:norm-one-grades}.

For $n\mid r$, use commutators between levels $1$ and $r-1$.
Neither level is divisible by $n$, since $n\geq3$.  Choose leading
coefficient $1$ at level $1$ and arbitrary coefficient $b$ at
level $r-1$.  The leading coefficient of their commutator is
$\sigma(b)-b$.  These fill
\[
 (\sigma-1)\bbF_{q^n}=\ker\Tr_{\bbF_{q^n}/\bbF_q}.
\]
Successive approximation in the complete filtration puts all of
$S_1$ in $\overline{[S_K,S_K]}$.  The reverse inclusion holds because
the residue quotient is abelian, proving the first assertion.

Finally, commutators of $\Pi$ with units of $E_K$ have residues
$\sigma(z)/z$, which fill $T_K$ by finite-field Hilbert~90.
Together with $S_1$ they are dense in $S_K$; all such commutators
have reduced norm one.  A unit of $\mathcal D$ acts trivially on
its commutative residue field, whereas $\Pi$ acts as $\sigma$.
Writing an arbitrary element as a unit times a power of $\Pi$
proves the assertion about conjugation.
\end{proof}

\subsection{A circle extension associated to the central quotient}

Write $I=\bbR/\bbZ$, additively, and fix an injective character
$\iota:C\hookrightarrow I$.  Give $L^1$ the discrete topology and set
\begin{equation}\label{quo:determinant-group}
 Q=\{(y,a)\in U_A\times L^1:\det y=a_A\},
 \qquad U_A=\prod_{v\in A}U(k_v).
\end{equation}
The determinant kernel is $S_A=\prod_{v\in A}S(k_v)$.
Thus $Q$ is locally compact and second countable, with countably many
open determinant fibers, each a translate of the compact group $S_A$.
The map
\[
 \rho:U(k)\longrightarrow Q,
 \qquad u\longmapsto(u_A,\det u)
\]
has dense image on every fiber by weak approximation for $S$.
It follows that $P_0$ is open and normal in $Q$ and that
\begin{equation}\label{quo:discrete-quotient}
 Q/P_0\simeq U(k)/V_0.
\end{equation}
For normality, first conjugate by the dense subgroup $\rho(U(k))$
and then pass to its closure; the quotient identification follows
from density and openness, with kernel $V_0$.

The central quotient in \eqref{quo:central-quotient} gives the discrete
central extension
\[
 1\longrightarrow C\longrightarrow U(k)/R'
   \longrightarrow U(k)/V_0\longrightarrow1.
\]
Pull it back along
$Q\to Q/P_0\simeq U(k)/V_0$, using \eqref{quo:discrete-quotient},
and push its kernel out by $\iota$.  We obtain a locally compact second-countable central
extension
\begin{equation}\label{quo:circle-extension}
 1\longrightarrow I\longrightarrow\widetilde Q
   \xrightarrow{p}Q\longrightarrow1.
\end{equation}
The pullback construction gives continuous local sections and an
abstract homomorphic lift
\[
 \ell:U(k)\longrightarrow\widetilde Q,
 \qquad p\ell=\rho.
\]
It also gives a canonical continuous homomorphic section
$c:P_0\to\widetilde Q$, obtained by using the identity in the
discrete group $U(k)/R'$.  For $u\in V_0$,
\begin{equation}\label{quo:canonical-lift}
 \ell(u)=c(u_A)\,\iota(uR').
\end{equation}
We regard values in $I$ as central factors in this multiplicative
notation for $\widetilde Q$.  Comparing $\ell$ with a continuous
splitting of $p$ will produce a character of $U(k)$.  Identity
\eqref{quo:canonical-lift} will ensure that its image on $S(k)$ is
finite and nonzero.

\begin{lemma}[Commuting lifts]\label{quo:commuting-lifts}
Suppose $q_1,q_2\in Q$ have local components in one common separable
field torus at each $v\in A$.  Any lifts of $q_1$ and $q_2$ to
$\widetilde Q$ commute.
\end{lemma}

\begin{proof}
Keep their two determinant coordinates fixed.  By
Proposition~\ref{quo:commuting}(i), their local components are limits
of commuting rational pairs with those exact determinants.  The
lifts under the homomorphism $\ell$ of each rational pair commute.
For commuting base elements, the commutator of lifts is independent
of the choices of lifts.  Continuous local sections of
\eqref{quo:circle-extension} therefore allow us to take the limit
of these commutators, giving the identity.
\end{proof}

The restriction of \eqref{quo:circle-extension} to $S_A$ has a
continuous homomorphic splitting.  Here is the metaplectic input with
its hypotheses made explicit.  Choose an isotropic place $v_0$ of
$S$ outside $A$, and pull the extension back to $S(\bbA^{v_0})$
by projection to $S_A$.  The lift $\ell|_{S(k)}$ splits this pullback
on rational points.  The metaplectic kernel away from $v_0$ is zero,
since $S$ is absolutely almost simple and simply connected and
$v_0$ is a nonarchimedean isotropic place
\cite[Theorem~3]{PRsurvey}; see also \cite{PRmetaplectic}.
The pullback extension consequently has zero measurable cohomology
class.  A measurable homomorphic section is continuous for these
locally compact second-countable groups.  Restrict it along the
subgroup $S_A\subset S(\bbA^{v_0})$ supported at $A$ to obtain
\begin{equation}\label{quo:local-section}
 j:S_A\longrightarrow\widetilde Q.
\end{equation}
This invocation involves only nonarchimedean places and therefore has
no omitted-archimedean-place qualification.

We must extend this splitting from $S_A$ to $Q$.  First we will make
$j(S_A)$ normal.  Its quotient will then be a central extension of
the abelian determinant group $L^1$; we will remove that extension's
commutators by a further change of $j$.

\subsection{Making the splitting equivariant}

For $q\in Q$ and any lift $\widetilde q$, define $d_q:S_A\to I$ by
\begin{equation}\label{quo:discrepancy}
 \widetilde q j(s)\widetilde q^{-1}
       =j(qsq^{-1})\,d_q(s).
\end{equation}
This is a continuous character, independent of the lift.  It is zero
for $q\in S_A$, and composition of conjugations gives
\[
 d_{q_1q_2}(s)=d_{q_1}(q_2sq_2^{-1})+d_{q_2}(s).
\]
Thus it is a character-valued cocycle factoring through
$Q/S_A=L^1$.  Lemma~\ref{quo:local-abelianization} identifies its
coefficient group as
\[
 \Hom_{\mathrm{cont}}(S_A,I)
       =\prod_{v\in A}\Hom(T_v,I),
 \qquad
 T_v=\ker N_{\bbF_{q_v^n}/\bbF_{q_v}}.
\]
On $T_v$, a determinant $a$ acts by $\sigma_v^{\ord_v(a_v)}$;
on its character group, the action in this cocycle is pullback,
$\lambda\mapsto\lambda\circ\sigma_v^{\ord_v(a_v)}$.

The $v$th component of $d_q$ is zero if $\ord_v(a_v)=0$.
Indeed, since $d_q$ depends only on the determinant, choose a
representative whose $v$th component is a unit in an unramified maximal subfield with reduced norm $a_v$; norms from that
field supply every unit.  It commutes with all Teichm\"uller lifts
of $T_v$, supported at $v$ in $S_A$.  At the other places choose
regular semisimple representatives of the specified determinant, so
that these components and the identity belong to separable field
tori.  Lemma~\ref{quo:commuting-lifts} then makes
\eqref{quo:discrepancy} zero on those Teichm\"uller lifts, which
determine the character.  The component and its action therefore
factor through the single valuation map $L^1\to\bbZ$ at $v$.
This map is onto by weak approximation in $L^1$.

Let $d_v$ be the discrepancy for valuation one.  It annihilates
$T_v^{\sigma_v}$: the Teichm\"uller lifts of this subgroup are
central scalars, and we may choose regular separable representatives
of the determinant and again apply Lemma~\ref{quo:commuting-lifts}.
For a finite abelian group $T$ with automorphism $\sigma$,
\[
 \im\bigl(\sigma^*-1:\Hom(T,I)\to\Hom(T,I)\bigr)
   =\{\lambda:\lambda|_{T^\sigma}=0\}.
\]
This is the dual of the kernel-image sequence for $\sigma-1$;
characters of a subgroup extend to $I$ because $I$ is divisible.
Choose a correcting character in each local factor, and let
$\lambda:S_A\to I$ be their sum.  For
$j_\lambda(s)=j(s)\lambda(s)$, direct substitution gives the discrepancy
\[
 d_q^\lambda(s)=d_q(s)+\lambda(s)-\lambda(qsq^{-1}).
\]
The characters may therefore be chosen to make this discrepancy zero
at the generator of each valuation group.  The cocycle identity then
makes every discrepancy zero.  Hence we may assume
\begin{equation}\label{quo:equivariant-section}
 \widetilde q j(s)\widetilde q^{-1}=j(qsq^{-1})
 \quad(q\in Q,\ s\in S_A).
\end{equation}
In particular, $j(S_A)$ is normal in $\widetilde Q$.  It is also
compact, hence closed, because $S_A$ is compact.

\subsection{Removing the determinant commutators}

The quotient by $j(S_A)$ is a central extension
\begin{equation}\label{quo:abelian-base-extension}
 1\longrightarrow I\longrightarrow\widetilde Q/j(S_A)
        \longrightarrow L^1\longrightarrow1.
\end{equation}
Commutators define an alternating biadditive pairing
$\omega:L^1\times L^1\to I$, with multiplicative arguments.
Lemma~\ref{quo:commuting-lifts} shows that
\begin{equation}\label{quo:common-norm-test}
 \omega(a,b)=0
\end{equation}
whenever, at every $v\in A$, both $a_v$ and $b_v$ are norms from one
common maximal separable subfield of $\mathcal D_v$.

We reduce this pairing to characters of local residue units.  Put
\[
 U_0=\{a\in L^1:\ord_v(a_v)=0\text{ for every }v\in A\}.
\]
Using unramified maximal subfields in \eqref{quo:common-norm-test}
gives $\omega(U_0,U_0)=0$.  The valuation map
$L^1\to\bbZ^A$ is onto by weak approximation.  For each $v\in A$,
choose a maximal totally ramified separable field $E_v/k_v$ inside
$\mathcal D_v$.  Its norm group is open: the separable norm map has
surjective trace differential, and hence has open image on local
points.  That norm group contains an element of valuation one.
Weak approximation now gives elements $\pi_i\in L^1$, $i\in A$,
whose valuation vectors are the standard basis vectors and all of
whose components at $v$ belong to $N_{E_v/k_v}(E_v^\times)$.
The common-norm test gives $\omega(\pi_i,\pi_j)=0$ for all $i,j$.
Every element of $L^1$ is a product of an element of $U_0$ and powers
of the $\pi_i$.  Only the pairings $\omega(\pi_v,u)$, $u\in U_0$,
remain to be considered.

For fixed $v$, this pairing vanishes whenever
$u_v\in N_{E_v/k_v}(E_v^\times)$: use $E_v$ at $v$ and unramified
maximal subfields at the other places, where both determinants are
units.  Put
\[
 H_v=\mathcal O_v^\times\cap N_{E_v/k_v}(E_v^\times).
\]
The group $H_v$ is open in the compact unit group, so its index is
finite.  Weak approximation makes the projection
$U_0\to\mathcal O_v^\times/H_v$ onto.  The vanishing just proved
says that its kernel is killed by $u\mapsto\omega(\pi_v,u)$.
Thus this pairing factors through that finite quotient.  Composing
the quotient character with the projection from $\mathcal O_v^\times$
defines a continuous character $\eta_v$ extending the pairing.
Since the norm group contains all $n$th powers, $n\eta_v=0$.
This argument uses openness of a separable norm group, not openness
of the $n$th-power subgroup.

We claim that $\eta_v$ kills the principal units.  Let $c\in k_v$
have valuation one.  Perturb the Eisenstein polynomial $X^n-c$, with
its constant coefficient fixed, to separable Eisenstein polynomials
$f_j$ tending to it.  If the characteristic divides $n$, one may take
$f_j=X^n+\varepsilon_jX-c$ with nonzero
$\varepsilon_j\to0$.  If it does not divide $n$, no perturbation is
necessary.  Let $\alpha_j$ be a root and $E'_j=k_v(\alpha_j)$.
These degree-$n$ separable fields embed in $\mathcal D_v$, and oddness
of $n$ gives
\[
 N_{E'_j/k_v}(\alpha_j)=c,
 \qquad
 N_{E'_j/k_v}(1+\alpha_j)\longrightarrow1+c.
\]
For each $j$, choose $\pi'_v\in L^1$ with the same valuation vector
as $\pi_v$, with $v$th component in $N_{E'_j/k_v}(E_j'{}^\times)$.
At the other places choose unit components.  Weak approximation and
openness of the norm groups permit this choice.  Since
$\pi'_v/\pi_v\in U_0$, their pairings with $U_0$ agree.
Choose elements of $U_0$ whose $v$th components belong to the same
norm group and tend to $N_{E'_j/k_v}(1+\alpha_j)$.
At the other places their components, and those of $\pi'_v$, are
norms from unramified maximal subfields.  The common-norm test and
continuity give
\[
 \eta_v\bigl(N_{E'_j/k_v}(1+\alpha_j)\bigr)=0.
\]
Letting $j$ tend to infinity gives $\eta_v(1+c)=0$.
These units generate $1+\mathfrak p_v$: if $\ord_v(d)\geq2$, then
\[
 1+d=\frac{(1+c)(1+d)}{1+c},
\]
and both numerator and denominator have the form $1+c'$ with
$\ord_v(c')=1$.  Thus $\eta_v$ is a character of
$\bbF_{q_v}^\times$, killed by $n$.

We can remove these residual characters without losing
\eqref{quo:equivariant-section}.  The permissible changes of $j$ are
characters of $S_A$ invariant under $Q$, whose $v$th factors are the
characters of $T_v/(\sigma_v-1)T_v$.  If a section is multiplied by
such a character $\lambda$, the commutator pairing changes by
$-\lambda$ evaluated on the ordinary local commutator of determinant
representatives.  To see this, write the commutator of any two lifts
as $j$ of their base commutator times the old value of $\omega$;
substitution of $j\lambda$ gives the asserted formula.

To compute the effect on $\eta_v$, represent the determinant
$(\pi_v)_v$ by a division prime multiplied by a unit of the fixed
unramified maximal subfield.  Such a unit adjusts the reduced norm to the specified value because
the ratio of the two norms is a unit.  It commutes with every other
unit of that subfield, so this adjustment does not affect their commutator.
For an unramified unit with residue $r$, the residue of that
commutator is therefore $\sigma_v(r)/r$.  Changing the
norm preimage $r$ alters it by $(\sigma_v-1)T_v$, so there is a
well-defined homomorphism
\begin{equation}\label{quo:residue-commutator}
 \bbF_{q_v}^\times\longrightarrow T_v/(\sigma_v-1)T_v,
 \qquad N(r)\longmapsto\sigma_v(r)/r.
\end{equation}
It is onto by finite-field Hilbert~90.  Its target has order
$\gcd(n,q_v-1)$: on the cyclic group $T_v$ the cokernel of
$\sigma_v-1$ has the same order as its kernel, and
\[
 T_v^{\sigma_v}=\{z\in\bbF_{q_v}^\times:z^n=1\}.
\]
The domain in \eqref{quo:residue-commutator} is cyclic.  Its dual
therefore identifies all characters killed by $n$ with characters
of that target.  Choose the invariant local character that cancels
$\eta_v$, and make these changes for every $v\in A$.
At other places the determinant representatives for a mixed pair can
be taken in unramified subfields, so no other local component is
introduced.  Unit-unit pairings and pairings of the chosen $\pi_i$
still vanish by \eqref{quo:common-norm-test}, which holds for every
equivariant choice of $j$.  All the generating pairings now vanish, so $\omega=0$.

The extension \eqref{quo:abelian-base-extension} is consequently
abelian.  The divisible group $I$ is injective as an abelian group,
so the extension splits, continuously since its base $L^1$ is
discrete.  Take in $\widetilde Q$ the inverse image of the image of
such a section.  This subgroup meets $I$ trivially and maps onto $Q$;
hence it maps isomorphically to $Q$.  The inverse map is continuous
on each open determinant fiber, where it is a translate of $j$.
We have constructed a continuous homomorphic splitting
\begin{equation}\label{quo:full-section}
 s:Q\longrightarrow\widetilde Q
\end{equation}
of the original extension \eqref{quo:circle-extension}.

\subsection{Extracting a finite character}

We finish the proof of Proposition~\ref{quo:dichotomy} by comparing
\eqref{quo:full-section} with the rational lift.  Define
$\chi_0:U(k)\to I$ by
\begin{equation}\label{quo:difference-character}
 \ell(u)=s(\rho(u))\,\chi_0(u).
\end{equation}
Both maps are homomorphisms lifting $\rho$, and their difference is
central, so $\chi_0$ is a character.

Its image on $S(k)$ is finite.  For $u\in R'$, Equation
\eqref{quo:canonical-lift} reads $\ell(u)=c(u_A)$.
Thus $\chi_0|_{R'}$ is the restriction of the continuous character
$p_0\mapsto c(p_0)s(p_0)^{-1}\in I$ on the profinite group $P_0$.  Such a character has finite
image: its compact image is a profinite quotient, while a closed
subgroup of the circle is either finite or the circle itself.
Since $R'$ has finite index in $S(k)$, the image $\chi_0(S(k))$ is
also finite.

This image is nonzero.  Choose a nonidentity element $\beta\in C$
and a representative $u\in V_0$.  The subgroup $R$ is dense in
$P_0$, so there are $r_j\in R$ with $(ur_j)_A\to1$.
All these elements still represent $\beta$ modulo $R'$.
If $\chi_0$ vanished on $S(k)$, Equations
\eqref{quo:canonical-lift} and \eqref{quo:difference-character} would
imply
\[
 c((ur_j)_A)\,\iota(\beta)=s(\rho(ur_j)).
\]
Both continuous sections tend to the identity.  This would give
$\iota(\beta)=0$, contradicting injectivity of $\iota$.

Finally, we make the character finite on all of $U(k)$ without
changing its restriction to $S(k)$.  Let $\mathbb F_L$ be the full
constant field of $L$.  The divisor map on $L^\times$ has kernel
$\mathbb F_L^\times$ and image in the free abelian divisor group.
Every subgroup of a free abelian group is free, and an extension
of a free abelian group splits.  Therefore
\[
 L^1=T\oplus\bigoplus_{j\in J}\bbZ a_j
\]
for a finite torsion group $T$ and a possibly infinite basis indexed
by $J$.  Choose $u_j\in U(k)$ with $\det u_j=a_j$ and a character
$\psi:L^1\to I$ satisfying
$\psi(a_j)=\chi_0(u_j)$ and $\psi|_T=0$.  Then
\[
 \chi=\chi_0-\psi\circ\det
\]
vanishes on the chosen $u_j$ and equals $\chi_0$ on $S(k)$.
If $m$ annihilates $\chi_0(S(k))$ and $t$ annihilates $T$, then
$mt$ annihilates $\chi(U(k))$: remove the free determinant part by
products of the $u_j$, and the $t$th power of what remains belongs
to $S(k)$.  The subgroup of $I$ annihilated by $mt$ is finite.
Thus $\chi$ has finite image and is nonzero on $S(k)$, proving the
first alternative.

All empty-product cases are included.  In particular, if $A$ is empty,
then $Q=L^1$ and $S_A=1$; Lemma~\ref{quo:commuting-lifts} makes the
circle extension abelian immediately, and the character extraction
above is unchanged.  This completes the proof of
Proposition~\ref{quo:dichotomy}.
